\xmpheader 5/{Odds and ends}
\chardef \\ = `\\ % let \\ denote a backslash
\footline{\hfil{\tenit -- \folio\ --}\hfil} %\footline provives a footer line.
% Here it's a centered, italicized page number.
\TeX\ knows how to hypernate words, but it isn't infallible.  If you are discussing the chemical
${\it 5}$-[p-(Flouro\-sul\-fonyl)ben\-zoyl]-l,%
$N^6$-ethe\-no\-adeno\-sine and \TeX\ complains to you about an ``overfull hbox'', try inserting some 
``discretionary hyphens''.  The notation `{\tt \\-}' tells \TeX\ about a dis\-cre\-tion\-ary hypen, 
that is, one that it might not have inserted otherwise.
\medskip
{\raggedright You can typeset text unjustified, i.e., with an uneven right margin.  In the old days,
before word processors were common, typewritten documents were unjustified beause there was no
convenient alternative.  Some people prefer text to be unjustified so that the spacing between words 
can be uniform.  Most books are set with justified margin, but not all.\par}

\proclaim Assertion 27.  There is an easy way to typeset the headings of assertions, theorems, etc.

Here's an example of how to typeset an itemized list two levels deep.  If you need more levels,
you'll have to program it yourself, alas.
\smallskip
\item {1.} The is the first item.
\item {2.} This is the second item.  It consists of two paragraphs.  We've
indented the second paragraph so that you can easily see where it starts.
\item {} \indent The second paragraph has three subitems underneath it.
\itemitem {(a)} This is the first subitem.
\itemitem {(b)} This is the second subitem.
\itemitem {(c)} This is the third subitem.
\item{$\bullet$} This is a strange-looking item because it's completely different from the others.
\smallskip
\leftline{Here's a left-justified line.$\Leftarrow$}
\rightline{$\Rightarrow$Here's a right-justified line.}
\centerline{$\Rightarrow$Here's a centered line.$\Leftarrow$}
% Don't try to use these commands within a paragraph.
